\فصل{مطالب تکمیلی}
\label{chap:results-supplement}

این پیوست، شواهد تکمیلی فصل~\ref{chap:results} را گردآوری می‌کند.
همهٔ شکل‌ها از خروجی‌های ذخیره‌شده تهیه شدند.
هیچ آموزش، برازش یا حل عددی تازه‌ای انجام نشد.

\قسمت{اعتبارسنجی وزن منتخب جریان}
\label{sec:appendix-checkpoint}

این آزمون با مجموعهٔ هفت‌نمونه‌ای اکتشافی فصل نتایج متفاوت است.
شناسهٔ نه‌نوسانی این بخش نیز به ورودی مستقل تعلق دارد.
میدان جریان در کانال مستطیلی بی‌بعد با ابعاد $5\times1$ بررسی شد.
مختصات، سرعت و فشار مطابق رابطهٔ~\eqref{eq:forward-flow-scales} بی‌بعد شدند.
چهار ورودی در شکل~\ref{fig:results-forward-flow-inlets} نشان داده شده‌اند.
این ورودی‌ها بتا، تخت، نامتقارن و دارای نه نوسان هستند.
عدد رینولدز هر حالت در شکل و جدول آمده است.

\begin{figure}[htbp]
\centering
\includegraphics[width=0.95\textwidth]{figs/chapter4/forward-flow-inlets.pdf}
\caption{پروفایل‌های بی‌بعد ورودی در ارزیابی تکرارشوندهٔ عملگر جریان.}
\label{fig:results-forward-flow-inlets}
\end{figure}

مرجع با روش اجزای محدود و شرایط مرزی یکسان محاسبه شد.
شبکه و مرجع در $7081$ نقطهٔ مشترک مقایسه شدند.
تراکم نقاط نزدیک ورودی بیشتر بود.
با ریزترشدن مش، تغییر نسبی سرعت کمتر از $0.010\%$ ماند.
تغییر نسبی فشار در همان بررسی کمتر از $0.20\%$ بود.
فشار با مرجع خروجی مشترک مقایسه شد.
حل‌های مرجع در انتخاب وزن نقش داشتند.
تابع هزینهٔ آموزش از پاسخ‌های اجزای محدود استفاده نکرد.
دو ورودی رزروشده در این ارزیابی بررسی نشدند.

شکل~\ref{fig:results-forward-flow-beta} میدان ورودی بتا را نشان می‌دهد.
شکل~\ref{fig:results-forward-flow-nine} به ورودی دارای نه نوسان اختصاص دارد.
هر ردیف یک کمیت بی‌بعد را در سه پنل نشان می‌دهد.
پنل‌ها میدان مرجع، پیش‌بینی و اختلاف علامت‌دار را نمایش می‌دهند.
اختلاف از کم‌کردن مرجع از پیش‌بینی به دست می‌آید.
در حالت نوسانی، اختلاف مؤلفهٔ عرضی نزدیک ورودی متمرکزتر است.

\begin{figure}[htbp]
\centering
\includegraphics[width=0.98\textwidth,height=0.72\textheight,keepaspectratio]{figs/chapter4/forward-flow-beta.pdf}
\caption{مقایسهٔ میدان پیش‌بینی‌شده و مرجع برای ورودی بتا در $\mathrm{Re}=30$.}
\label{fig:results-forward-flow-beta}
\end{figure}

\begin{figure}[htbp]
\centering
\includegraphics[width=0.98\textwidth,height=0.72\textheight,keepaspectratio]{figs/chapter4/forward-flow-nine.pdf}
\caption{مقایسهٔ میدان پیش‌بینی‌شده و مرجع برای ورودی نوسانی در $\mathrm{Re}=55$.}
\label{fig:results-forward-flow-nine}
\end{figure}

جدول~\ref{tab:results-forward-flow-errors} خطاهای وزن منتخب را نشان می‌دهد.
این وزن در گام $94000$ ذخیره شد.
خطاهای نسبی مطابق رابطهٔ~\eqref{eq:validation-relative-l2} و برحسب درصد گزارش شده‌اند.
معیارهای ریشهٔ میانگین مربعات مطابق رابطهٔ~\eqref{eq:validation-forward-flow-rms} هستند.

\begin{table}[htbp]
\centering
\small
\setlength{\tabcolsep}{4pt}
\caption{خطاهای بی‌بعد میدان و باقیماندهٔ مومنتوم برای وزن منتخب.}
\label{tab:results-forward-flow-errors}
\begin{tabular}{@{}lcccccc@{}}
\hline
ورودی & $\mathrm{Re}$ & $E_2(\mathbf U^*)$ & $E_2(v^*)$ &
$E_{v,\mathrm{rms}}^*$ & $E_2(p^*)$ & $E_{\mathrm{mom,rms}}^*$ \\
\hline
بتا & $30$ & $0.906\%$ & $9.261\%$ & $0.004159$ & $3.082\%$ & $0.2759$ \\
تخت & $70$ & $3.002\%$ & $39.781\%$ & $0.010812$ & $10.974\%$ & $0.5170$ \\
نامتقارن & $100$ & $4.649\%$ & $37.160\%$ & $0.008145$ & $17.845\%$ & $0.1509$ \\
نه نوسان & $55$ & $6.604\%$ & $89.199\%$ & $0.024466$ & $8.210\%$ & $2.0929$ \\
\hline
\end{tabular}
\end{table}

ورودی بتا کمترین خطای بردار سرعت و فشار را داشت.
در حالت نوسانی، خطای نسبی مؤلفهٔ عرضی به $89.199\%$ رسید.
ریشهٔ میانگین مربعات همان مؤلفه برابر $0.024466$ بود.
بنابراین دقت مؤلفهٔ عرضی باید جدا از خطای بردار سرعت سنجیده شود.
بیشترین خطای نسبی فشار در حالت نامتقارن ثبت شد.
در این مجموعه، شکل ورودی و عدد رینولدز هم‌زمان تغییر کردند.

حد پذیرش خطای نسبی سرعت $5\%$ و فشار $10\%$ بود.
خطای سرعت در حالت نوسانی از حد تعیین‌شده گذشت.
خطای فشار در حالت‌های تخت و نامتقارن از حد گذشت.
باقیماندهٔ مومنتوم در هر چهار حالت از آستانهٔ $0.05$ بیشتر بود.
معیارهای پیوستگی، شرط خروجی، دبی و خطای ورودی پذیرفته شدند.
در مجموع، وزن منتخب سه معیار از هفت معیار را برآورده نکرد.
وزن گام $100000$ نیز همان سه معیار را برآورده نکرد.
این ارزیابی، مبنای گزارش محدودیت میدان در بازسازی معکوس است.


\FloatBarrier
\قسمت{شکل‌های تکمیلی پیش‌بینی مستقیم}
\label{sec:appendix-forward}

\زیرقسمت{ورودی‌های اکتشافی جریان}
چهار نمونهٔ باقی‌مانده در شکل‌های بعدی ارائه شده‌اند.
تعریف نرم و وزن نقاط با جدول~\ref{tab:results-forward-flow-exploratory} یکسان است.
نسخهٔ اصلی شکل‌های جریان نیز در مخزن پژوهشی در دسترس است.
\begin{latin}
\small
\href{https://github.com/mohrezsayeh90-ai/inverse_fromtrajectory_to_FlowReconstruction/blob/92897a6962e537a4bea88042609be102095db037/results/diverse_profiles_2026_11_07/all_new_figures.pdf}{Original seven-case flow and timed inverse gallery}
\end{latin}
خطای نسبی مؤلفهٔ عرضی به کوچک‌بودن مرجع حساس است.
ریشهٔ میانگین مربعات آن، معیار مکمل با مقیاس ثابت است.

شکل~\ref{fig:appendix-flow-waves3_r20} نمونهٔ سه‌دوره‌ای را نشان می‌دهد.

\begin{figure}[htbp]
\centering
\includegraphics[width=\textwidth]{figs/chapter4/forward-flow-chain-waves3_r20.pdf}
\caption{میدان، خطا و برش‌های سرعت برای ورودی سه‌دوره‌ای.}
\label{fig:appendix-flow-waves3_r20}
\end{figure}

\FloatBarrier

شکل~\ref{fig:appendix-flow-lobes4_r20} نمونهٔ چهاردوره‌ای را نشان می‌دهد.

\begin{figure}[htbp]
\centering
\includegraphics[width=\textwidth]{figs/chapter4/forward-flow-chain-lobes4_r20.pdf}
\caption{میدان، خطا و برش‌های سرعت برای ورودی چهاردوره‌ای.}
\label{fig:appendix-flow-lobes4_r20}
\end{figure}

\FloatBarrier

شکل~\ref{fig:appendix-flow-mixed} نمونهٔ ترکیبی را نشان می‌دهد.

\begin{figure}[htbp]
\centering
\includegraphics[width=\textwidth]{figs/chapter4/forward-flow-chain-mixed.pdf}
\caption{میدان، خطا و برش‌های سرعت برای ورودی ترکیبی.}
\label{fig:appendix-flow-mixed}
\end{figure}

\FloatBarrier

شکل~\ref{fig:appendix-flow-shoulder} نمونهٔ شانه‌دار را نشان می‌دهد.

\begin{figure}[htbp]
\centering
\includegraphics[width=\textwidth]{figs/chapter4/forward-flow-chain-shoulder.pdf}
\caption{میدان، خطا و برش‌های سرعت برای ورودی شانه‌دار.}
\label{fig:appendix-flow-shoulder}
\end{figure}

\FloatBarrier

\زیرقسمت{فرکانس متغیر و رهاسازی پایین‌دست مغناطیسی}
شکل~\ref{fig:appendix-magnetic-09} نمونهٔ فرکانس متغیر دو منبع را نشان می‌دهد.
خطای همین نمونه در شکل~\ref{fig:results-forward-magnetic-errors} نیز آمده است.

\begin{figure}[htbp]
\centering
\includegraphics[width=\textwidth]{figs/chapter4/forward-magnetic-chain-09.pdf}
\caption{تحریک دو منبع با فرکانس متغیر و پاسخ ذره.}
\label{fig:appendix-magnetic-09}
\end{figure}

\FloatBarrier
شکل~\ref{fig:appendix-magnetic-11} پاسخ رهاسازی پایین‌دست را نشان می‌دهد.

\begin{figure}[htbp]
\centering
\includegraphics[width=\textwidth]{figs/chapter4/forward-magnetic-chain-11.pdf}
\caption{تغییر منبع و پاسخ ذره با رهاسازی پایین‌دست.}
\label{fig:appendix-magnetic-11}
\end{figure}

نسخهٔ اصلی دوازده زنجیره در مخزن پژوهشی در دسترس است.
\begin{latin}
\small
\href{https://github.com/mohrezsayeh90-ai/Master_Magnetic/blob/639906f53a11934b3058fbbdb246e34a1dd001fd/results/trained_presentation/trained_cases_panels_1_to_5.pdf}{Original 12-case magnetic gallery}
\end{latin}

\FloatBarrier
\قسمت{مقایسه‌های تکمیلی بازسازی جریان}
\label{sec:appendix-inverse}

\زیرقسمت{بازپخش زمان‌دار و مقایسه با میدان مرجع}
نمونهٔ گوشه‌ای الف با نمونهٔ ب در متن اصلی متفاوت است.
شکل~\ref{fig:results-inverse-timed-replay} بازپخش نمونهٔ گوشه‌ای الف را نشان می‌دهد.
در بازپخش، معادلهٔ حرکت از شرایط اولیه دوباره انتگرال‌گیری شد.
سرعت سیال در مکان جاری ذره خوانده شد.
خطای مکانی در زمان‌های مشاهده محاسبه شد.
پروفایل ورودی تخمینی نیز با پروفایل مرجع مقایسه شده است.

\begin{figure}[htbp]
\centering
\includegraphics[width=0.97\textwidth,height=0.68\textheight,keepaspectratio]{figs/chapter4/inverse-timed-replay.pdf}
\caption{مقایسهٔ مسیرها و پروفایل ورودی در بازسازی زمان‌دار نمونهٔ گوشه‌ای الف.}
\label{fig:results-inverse-timed-replay}
\end{figure}

شکل~\ref{fig:results-inverse-timed-fields} سه مؤلفهٔ میدان همان نمونه را نشان می‌دهد.
هر ردیف شامل میدان مرجع، میدان تخمینی و اختلاف علامت‌دار است.
اختلاف از کم‌کردن مرجع از تخمین به دست می‌آید.
اختلاف در بخش‌هایی از کانال، دور از مسیرها نیز دیده می‌شود.

\begin{figure}[htbp]
\centering
\includegraphics[width=0.98\textwidth,height=0.72\textheight,keepaspectratio]{figs/chapter4/inverse-timed-fields.pdf}
\caption{مقایسهٔ میدان بازسازی‌شده و مرجع اجزای محدود در نمونهٔ گوشه‌ای الف.}
\label{fig:results-inverse-timed-fields}
\end{figure}

جدول~\ref{tab:results-inverse-timed} خطاهای مسیر، پروفایل و میدان را گزارش می‌کند.
معیار مسیر از رابطهٔ~\eqref{eq:validation-inverse-timed-path} به دست می‌آید.
خطاهای نسبی میدان مطابق رابطهٔ~\eqref{eq:validation-relative-l2} محاسبه شدند.
هر شش بازپخش به پایان بازهٔ مشاهده رسیدند.
بهینه‌سازی در شروع‌های منتخب با رسیدن به سقف تکرار پایان یافت.

\begin{table}[htbp]
\centering
\small
\setlength{\tabcolsep}{4pt}
\caption{خطاهای بی‌بعد بازپخش و مقایسهٔ بازسازی زمان‌دار با مرجع.}
\label{tab:results-inverse-timed}
\begin{tabular}{@{}lcccccc@{}}
\hline
نمونه & $10^4 E_{1,\mathrm{rms}}^*$ & $10^4 E_{2,\mathrm{rms}}^*$ & $A$ &
$E_2(f_{\mathrm{ref}}^*)$ & $E_2(\mathbf U_{\mathrm{ref}}^*)$ & $E_2(p_{\mathrm{ref}}^*)$ \\
\hline
گوشه‌ای ب & $4.443$ & $3.768$ & $1.00194$ & $32.423\%$ & $19.505\%$ & $20.795\%$ \\
گوشه‌ای الف & $6.676$ & $4.821$ & $1.00236$ & $43.153\%$ & $26.913\%$ & $32.999\%$ \\
مایل & $18.984$ & $10.001$ & $1.01047$ & $40.213\%$ & $21.217\%$ & $42.804\%$ \\
\hline
\end{tabular}
\end{table}

در جدول، $f_{\mathrm{ref}}^*$ پروفایل ورودی در مقیاس سرعت مرجع است.
در نمونهٔ گوشه‌ای الف، خطای بازپخش هر دو ذره کمتر از $7\times10^{-4}$ بود.
خطای نسبی پروفایل همان نمونه به $43.153\%$ رسید.
خطای نسبی میدان سرعت نیز $26.913\%$ بود.
فشار از پاسخ اپراتور جریان به دست آمد.
این مقایسه نشان می‌دهد بازتولید دو مسیر برای تأیید کل میدان کافی نیست.
تفسیر میدان به محدودیت‌های عملگر مستقیم در بخش~\ref{sec:results-forward-flow} وابسته است.


شکل~\ref{fig:appendix-timed-a} زنجیرهٔ زمان‌دار نمونهٔ گوشه‌ای الف را نشان می‌دهد.

\begin{figure}[htbp]
\centering
\includegraphics[width=\textwidth]{figs/chapter4/inverse-timed-chain-a.pdf}
\caption{مسیرها، میدان تخمینی و خطای بازپخش زمان‌دار نمونهٔ گوشه‌ای الف.}
\label{fig:appendix-timed-a}
\end{figure}

\FloatBarrier

شکل~\ref{fig:appendix-timed-oblique} زنجیرهٔ زمان‌دار نمونهٔ مایل را نشان می‌دهد.

\begin{figure}[htbp]
\centering
\includegraphics[width=\textwidth]{figs/chapter4/inverse-timed-chain-oblique.pdf}
\caption{مسیرها، میدان تخمینی و خطای بازپخش زمان‌دار نمونهٔ مایل.}
\label{fig:appendix-timed-oblique}
\end{figure}

\FloatBarrier

\زیرقسمت{هندسه و ساعت تخمینی نمونه‌های گوشه‌ای}
\begin{latin}
\small
\href{https://github.com/mohrezsayeh90-ai/inverse_fromtrajectory_to_FlowReconstruction/blob/92897a6962e537a4bea88042609be102095db037/flow_reconstruction_no_time/chain_revision_xt_yt/chains_only.pdf}{Original three-case geometry-only gallery}
\end{latin}
نمونهٔ ب، همان دادهٔ شکل~\ref{fig:results-inverse-timed-chain-b} را بدون زمان بررسی می‌کند.
شکل~\ref{fig:appendix-geometry-b} این ارتباط را نشان می‌دهد.

\begin{figure}[htbp]
\centering
\includegraphics[width=\textwidth]{figs/chapter4/inverse-geometry-chain-b.pdf}
\caption{بازسازی هندسی و ساعت تخمینی در نمونهٔ گوشه‌ای ب.}
\label{fig:appendix-geometry-b}
\end{figure}

\FloatBarrier
شکل~\ref{fig:appendix-geometry-a} به نمونهٔ گوشه‌ای الف تعلق دارد.

\begin{figure}[htbp]
\centering
\includegraphics[width=\textwidth]{figs/chapter4/inverse-geometry-chain-a.pdf}
\caption{بازسازی هندسی و ساعت تخمینی در نمونهٔ گوشه‌ای الف.}
\label{fig:appendix-geometry-a}
\end{figure}

\FloatBarrier
جدول~\ref{tab:results-inverse-geometry} هر شش مسیر و دو معیار را مقایسه می‌کند.
\begin{table}[htbp]
\centering
\small
\setlength{\tabcolsep}{5pt}
\caption{خطاهای هندسی بی‌بعد دو بازپخش با مقیاس مشترک.}
\label{tab:results-inverse-geometry}
\begin{tabular}{@{}lccccc@{}}
\hline
 & & \multicolumn{2}{c}{$10^4 E_{\mathrm{sym}}^*$}
     & \multicolumn{2}{c}{$10^4 E_{\mathrm{arc}}^*$} \\
نمونه & ذره & زمان‌دار & بدون زمان & زمان‌دار & بدون زمان \\
\hline
گوشه‌ای ب & $1$ & $2.713$ & $2.130$ & $6.467$ & $2.134$ \\
گوشه‌ای ب & $2$ & $2.059$ & $4.030$ & $8.432$ & $19.779$ \\
گوشه‌ای الف & $1$ & $3.388$ & $3.149$ & $7.773$ & $3.152$ \\
گوشه‌ای الف & $2$ & $2.219$ & $4.238$ & $10.850$ & $21.478$ \\
مایل & $1$ & $13.067$ & $9.523$ & $14.242$ & $10.932$ \\
مایل & $2$ & $8.795$ & $2.264$ & $15.992$ & $4.228$ \\
\hline
\end{tabular}
\end{table}
\begin{figure}[htbp]
\centering
\includegraphics[width=0.97\textwidth,height=0.75\textheight,keepaspectratio]{figs/chapter4/inverse-geometry-comparison.pdf}
\caption{مقایسهٔ هندسهٔ بازپخش و پروفایل‌های تخمینی در دو روش.}
\label{fig:results-inverse-geometry}
\end{figure}
\FloatBarrier
\زیرقسمت{چیدمان مطالعهٔ تعداد مسیرها}
شکل~\ref{fig:inverse-count-design} ترتیب مسیرهای افزوده‌شده را مشخص می‌کند.
این مطالعه از آزمایش بازپخش نمونهٔ ب مستقل است.
\begin{figure}[htbp]
\centering
\includegraphics[width=0.92\textwidth]{figs/chapter4/inverse-particle-count-design.pdf}
\caption{ترتیب شش مسیر مرجع در کانال بی‌بعد برای ورودی نه‌نوسانی.}
\label{fig:inverse-count-design}
\end{figure}
\FloatBarrier
\قسمت{جزئیات بازیابی تحریک مغناطیسی}
\label{sec:appendix-magnetic-inverse}

جدول~\ref{tab:results-master10-inverse} آزمون‌های حساسیت و نویز را نیز شامل می‌شود.
شکل‌های این بخش، تک‌منبعی و دومنبعی را جداگانه مقایسه می‌کنند.
موقعیت‌های زمان‌دار مصنوعی برای بازیابی دو دامنهٔ منبع استفاده شدند.
جریان حامل، هندسه و پارامترهای سیال معلوم و ثابت بودند.
مدل بیش‌میرا به سرعت اولیهٔ مستقل نیاز ندارد.
شرط اولیه از اولین موقعیت دقیق گرفته شد.
وزن‌های آموزش‌دیدهٔ مدل دهم در تمام برازش‌ها ثابت ماندند.
روش از هدف رابطهٔ~\eqref{eq:inverse-master10-proposed-objective} استفاده می‌کند.

در آزمون پایه، هفت گره روی بازهٔ $[0,1.5]$ قرار داشتند.
بین گره‌ها از درون‌یابی خطی استفاده شد.
پنجرهٔ پایه برابر $0.0125$ و محاسبات از نوع \lr{float64} بود.
معیار انتخاب شروع، مجموع هزینهٔ داده و جریمه‌های اعلام‌شده بود.
شروع‌ها غیرصفر و مستقل از تحریک حقیقی انتخاب شدند.
تحریک مرجع فقط در تولید مشاهده و ارزیابی استفاده شد.
بازپخش مستقل، ضرایب خام تخمینی را دریافت کرد.

در آزمون $A$، فقط منبع چپِ مثبت مجهول بود.
در آزمون $B$، هر دو منبع بدون پیشین علامت تخمین زده شدند.
هر دو آزمون از موقعیت اولیهٔ $(X_0,Z_0)=(0.12,-1)$ آغاز شدند.
هر مسیر پایه $121$ مشاهده داشت.
ارزیابی مستقل روی $601$ زمان مشترک انجام شد.
آزمون $C$ از حل عددیِ تحریک نرم تغییر منبع استفاده کرد.
افق زمانی این آزمون $2.25$ بود.

شکل~\ref{fig:results-master10-controls} تحریک حقیقی و تخمینی دو آزمون را نشان می‌دهد.
برای مقایسه، یک علامت مشترک برای هر دو منبع انتخاب شد.
این علامت در همهٔ زمان‌ها ثابت ماند.
معیار مقایسه مطابق رابطهٔ~\eqref{eq:validation-master10-proposed-recovery} است.

\begin{figure}[htbp]
\centering
\includegraphics[width=0.98\textwidth]{figs/chapter4/inverse-master10-controls.pdf}
\caption{تحریک مرجع و تخمین هم‌تراز در آزمون‌های تک‌منبعی و دومنبعی.}
\label{fig:results-master10-controls}
\end{figure}

شکل~\ref{fig:results-master10-replay} پیش‌بینی شبکه و بازپخش مستقل را مقایسه می‌کند.
بازپخش با حل‌گر \lr{DOP853} و معادلهٔ اصلی انجام شد.
مقیاس‌های عرضی و محوری، طول‌های مرجع متفاوت دارند.
قاب مسیر بزرگ‌نمایی شده و مرز آن دیوارهٔ لوله نیست.
خطای هر مؤلفه جداگانه و با مقیاس ثابت گزارش می‌شود.

\begin{figure}[htbp]
\centering
\includegraphics[width=0.98\textwidth]{figs/chapter4/inverse-master10-replay.pdf}
\caption{مسیر تخمینی و بازپخش مستقل همراه خطاهای بی‌بعد دو مؤلفه.}
\label{fig:results-master10-replay}
\end{figure}

جدول~\ref{tab:results-master10-inverse} بیشینهٔ خطاهای بازپخش نسبت به مرجع پاک را نشان می‌دهد.
خطای تحریک، نرم نسبی مشترک دو منبع پس از هم‌ترازی است.
وضعیت جدول فقط به توقف بهینه‌ساز منتخب اشاره دارد.
تمام بازپخش‌های مستقلِ جدول کامل و داخل دامنه بودند.
سقف هر شروع $90$ تکرار بود.
در حالت پایه سه شروع و در مقایسه‌ها دو شروع اجرا شد.

\begin{table}[htbp]
\centering
\small
\setlength{\tabcolsep}{3pt}
\caption{خطای بازپخش، خطای تحریک و وضعیت بهینه‌ساز در آزمایش‌های معکوس.}
\label{tab:results-master10-inverse}
\begin{tabular}{@{}llcccc@{}}
\hline
آزمون & تنظیم & $10^5 E_{X,\infty}$ & $10^5 E_{Z,\infty}$ & $E_a(\%)$ & توقف\\
\hline
\lr{A} & منبع چپ مثبت & $6.315$ & $1.144$ & $0.37$ & توقف خطی \\
\lr{B} & دو منبع؛ علامت نامعلوم & $46.592$ & $3.251$ & $104.65$ & سقف تکرار \\
\lr{Bk} & سیزده گره & $76.918$ & $3.165$ & $118.75$ & سقف تکرار \\
\lr{Bh} & پنجرهٔ نصف & $46.480$ & $3.006$ & $104.82$ & سقف تکرار \\
\lr{Bw} & مقیاس محوری دوبرابر & $46.842$ & $2.922$ & $104.87$ & سقف تکرار \\
\lr{Bl} & زبری ده‌برابر & $47.643$ & $3.261$ & $105.27$ & سقف تکرار \\
\lr{C} & تحریک نرم بایگانی‌شده & $89.340$ & $8.349$ & $102.44$ & سقف تکرار \\
\lr{Ds} & سی‌ویک مشاهده & $46.365$ & $2.760$ & $103.91$ & سقف تکرار \\
\lr{D1} & نویز $0.05$؛ بذر $1103$ & $54.330$ & $1.863$ & $100.96$ & سقف تکرار \\
\lr{D2} & نویز $0.05$؛ بذر $2907$ & $53.364$ & $1.667$ & $104.34$ & سقف تکرار \\
\lr{D3} & نویز $0.2$؛ بذر $1103$ & $184.551$ & $2.863$ & $98.56$ & سقف تکرار \\
\lr{D4} & نویز $0.2$؛ بذر $2907$ & $71.093$ & $1.193$ & $110.21$ & سقف تکرار \\
\hline
\end{tabular}
\end{table}

در دادهٔ کم‌نمونه، هر چهارمین مشاهده و دو انتها حفظ شدند.
نویز جدول، انحراف معیار فیزیکی برحسب میکرومتر برای هر مؤلفه است.
نویز هر مؤلفه با طول مرجع همان مؤلفه بی‌بعد شد.
شرط اولیه در تمام آزمون‌های نویزی دقیق باقی ماند.
مقیاس برازش، بیشینهٔ نویز معلوم و کف خطای عملگر بود.

خطای تحریک تک‌منبعی حدود $0.374\%$ بود.
در آزمون دومنبعی پایه، این خطا به $104.65\%$ رسید.
برازش منتخب تک‌منبعی با توقف جست‌وجوی خطی پایان یافت.
دو شروع دیگر تک‌منبعی همگرا شدند.
این تفاوت، دقت مسیر را از دقت تحریک متمایز می‌کند.
توقف در سقف تکرار، همگرایی یا کمینهٔ سراسری را تأیید نمی‌کند.
نتایج ناموفقِ بازیابی نیز در جدول حفظ شده‌اند.
این نتایج به بودجه، پیشین‌ها و نمایش خطی همین آزمایش محدودند.
ارزیابی مستقیم قبلی با این برازش‌ها تغییر نمی‌کند.


\FloatBarrier
\قسمت{بازتاب برنامه و حساسیت رهاسازی}
\label{sec:appendix-formation}

\زیرقسمت{منابع پشتیبان هندسه و پوشش}
چیدمان ایستگاه‌ها، پوشش هدف و برش‌های مرکز در منابع پشتیبان ثبت‌اند.
پوشش هندسی با نرخ برخورد به هدف یکسان نیست.
\begin{latin}
\small
\href{https://github.com/mohrezsayeh90-ai/Main_Proposal/tree/9690560dd817fc3bea96007fd4c200a1141db712/New_Main_Results/runs/formation_20261112_cpu/figures}{Formation support: F03\_ports, F05\_occupancy, F06\_slices}
\end{latin}



\زیرقسمت{بازتاب یک برنامهٔ صدذره‌ای ذخیره‌شده}
شکل~\ref{fig:appendix-formation-reflection} پاسخ بازتابی را با تبدیل هندسی مقایسه می‌کند.
بیشینهٔ باقیماندهٔ مرجع تحلیلی برابر $6.80436\times10^{-16}$ بود.
برای میدان شبکهٔ ثابت، این مقدار $1.093\times10^{-4}$ بود.
مقایسهٔ بهینه‌سازی دسترس‌پذیری دیواره‌ها همچنان انجام نشده است.

\begin{figure}[htbp]
\centering
\includegraphics[width=\textwidth]{figs/chapter4/formation-reflection.pdf}
\caption{مقایسهٔ بازتاب هندسی و بازپخش مستقل یک برنامهٔ صدذره‌ای.}
\label{fig:appendix-formation-reflection}
\end{figure}

\FloatBarrier
\زیرقسمت{حساسیت عددی برنامه‌های منتخب}
شکل~\ref{fig:appendix-formation-sensitivity} دو نمایندهٔ کم‌حاشیه و پُرحاشیه را مقایسه می‌کند.
تعداد ذرات این نماینده‌ها حداکثر $24$ بود.
برای هر وضعیت تکمیل‌شده، $64$ اغتشاش بررسی شد.
آستانهٔ پذیرش خطای مرجع $0.004$ بود.
زمان هدف اسمی در هر نماینده ثابت ماند.
پذیرش نامشروط و مشروط به قیود، جداگانه نمایش داده شدند.
با نبود نمونهٔ مجاز، کسر پذیرش مشروط تعریف نمی‌شود.
این حساسیت عددی، معیار قابلیت‌اعتماد سخت‌افزار نیست.

\begin{figure}[htbp]
\centering
\includegraphics[width=\textwidth]{figs/chapter4/formation-release-sensitivity.pdf}
\caption{حساسیت عددی برنامه‌های منتخب به اغتشاش رهاسازی و زمان‌بندی.}
\label{fig:appendix-formation-sensitivity}
\end{figure}
